\documentclass[12pt]{article}
\usepackage{amsmath,amssymb}

\begin{document}
\section*{{2022 AIME I Problems/Problem 11}}
Source: AoPS Wiki

\subsection*{Solution}
Let's redraw the diagram, but extend some helpful lines. [asy] size(10cm); pair A,B,C,D,EE,F,P,Q,O; A=(0,0); EE = (24,15); F = (30,0); O = (10.5,7.5); label("$A$", A, SW); B=(6,15); label("$B$", B, NW); C=(30,15); label("$C$", C, NE); D=(24,0); label("$D$", D, SE); P=(5.2,2.6); label("$P$", (5.8,2.6), N); Q=(18.3,9.1); label("$Q$", (18.1,9.7), W); draw(A--B--C--D--cycle); draw(C--A); draw(Circle((10.95,7.45), 7.45)); dot(A^^B^^C^^D^^P^^Q); dot(O); label("$O$",O,W); draw((10.5,15)--(10.5,0)); draw(D--(24,15),dashed); draw(C--(30,0),dashed); draw(D--(30,0)); dot(EE); dot(F); label("$3$", midpoint(A--P), S); label("$9$", midpoint(P--Q), S); label("$16$", midpoint(Q--C), S); label("$x$", (5.5,13.75), W); label("$20$", (20.25,15), N); label("$6$", (5.25,0), S); label("$6$", (1.5,3.75), W); label("$x$", (8.25,15),N); label("$14+x$", (17.25,0), S); label("$6-x$", (27,15), N); label("$6+x$", (27,7.5), S); label("$6\sqrt{3}$", (30,7.5), E); label("$T_1$", (10.5,15), N); label("$T_2$", (10.5,0), S); label("$T_3$", (4.5,11.25), W); label("$E$", EE, N); label("$F$", F, S); [/asy] We obviously see that we must use power of a point since they've given us lengths in a circle and there are intersection points. Let $T_1, T_2, T_3$ be our tangents from the circle to the parallelogram. By the secant power of a point, the power of $A = 3 \cdot (3+9) = 36$ . Then $AT_2 = AT_3 = \sqrt{36} = 6$ . Similarly, the power of $C = 16 \cdot (16+9) = 400$ and $CT_1 = \sqrt{400} = 20$ . We let $BT_3 = BT_1 = x$ and label the diagram accordingly. Notice that because $BC = AD, 20+x = 6+DT_2 \implies DT_2 = 14+x$ . Let $O$ be the center of the circle. Since $OT_1$ and $OT_2$ intersect $BC$ and $AD$ , respectively, at right angles, we have $T_2T_1CD$ is a right-angled trapezoid and more importantly, the diameter of the circle is the height of the triangle. Therefore, we can drop an altitude from $D$ to $BC$ and $C$ to $AD$ , and both are equal to $2r$ . Since $T_1E = T_2D$ , $20 - CE = 14+x \implies CE = 6-x$ . Since $CE = DF, DF = 6-x$ and $AF = 6+14+x+6-x = 26$ . We can now use Pythagorean theorem on $\triangle ACF$ ; we have $26^2 + (2r)^2 = (3+9+16)^2 \implies 4r^2 = 784-676 \implies 4r^2 = 108 \implies 2r = 6\sqrt{3}$ and $r^2 = 27$ . We know that $CD = 6+x$ because $ABCD$ is a parallelogram. Using Pythagorean theorem on $\triangle CDF$ , $(6+x)^2 = (6-x)^2 + 108 \implies (6+x)^2-(6-x)^2 = 108 \implies 12 \cdot 2x = 108 \implies 2x = 9 \implies x = \frac{9}{2}$ . Therefore, base $BC = 20 + \frac{9}{2} = \frac{49}{2}$ . Thus the area of the parallelogram is the base times the height, which is $\frac{49}{2} \cdot 6\sqrt{3} = 147\sqrt{3}$ and the answer is $\boxed{150}$ ~KingRavi

\end{document}
